\documentclass[../candidatura.tex]{subfiles}
\graphicspath{{\subfix{../imagenes/}}}
\begin{document}

\section{Carcasas}

\avisoesbozo

\subsection{Carcasas paramétricas}

En mi experiencia de docente en diseño en Chile, veo que gran parte del trabajo de los estudiantes es hecho en software donde se privilegia el trabajo con el paradigma de hacer clic, de forma visual, no de forma textual ni paramétrica.

Se proponen formas paramétricas de desarrollar carcasas, donde en vez de pensar en casos puntuales, se sistematizan y publican herramientas para calcular y componer carcasas que permitan construir interfaces entre usuario y máquina, donde se vuelvan usables las perillas, los botones y los nodos de interconexión.

% \subsection{Carcasas en la obra de Critter and Guitari}

% \subsection{Carcasas en la obra de BASTL Instruments}

% \subsection{Carcasas en la obra de Corazón de Robota}

\subsection{Carcasas en la obra de Matías Serrano}

La obra de los últimos 15 años del artista Matías Serrano tiene un fuerte énfasis en circuitos, placas, instrumentos, interfaces. Las carcasas en sus obras muchas veces no están presentes, prefiriendo la estrategia y estética de circuitos expuestos, placas con potenciómetros expuestos, y que tanto documentación textual, como poemas y gráficas estén serigrafiadas directamente en la placa de circuito impreso.


\newpage

\end{document}
